% !TeX root = ../../Book1.tex
% 纲要标题：### 11.2 Euclid 整环
% 纲要位置：计划纲要.md，第 422 行。

\section{Euclid 整环}
\label{b1:sec:1102}

% 纲要内容（仅作写作计划，尚非教材正文）：
%
% * 除法算法与 Euclid 算法
% * 整数环及一元多项式环
% * Gauss 整数的例子
%

整数中求最大公因子的辗转相除法，依靠的是余数不断变小。推广这个方法时，真正需要的不是元素本身有大小，而是给非零余数配一个不能无限下降的非负整数。本节把这一要求写成公理，并在整数、多项式和复平面格点中实际执行除法。

\subsection{除法算法与 Euclid 算法}
\label{b1:subsec:1102:01}

\begin{definition}\label{b1:def:euclidean-domain}
设 \(R\) 是交换整环。若存在函数
\[
\delta:R\setminus\{0\}\longrightarrow\mathbb N
\]
使任意 \(a,b\in R\)、\(b\ne0\) 都可写成 \(a=bq+r\)，其中 \(q,r\in R\)，并且 \(r=0\) 或 \(\delta(r)<\delta(b)\)，则称 \(R\) 为\term[euclid-zhenghuan]{Euclid 整环}[Euclidean domain]，称 \(\delta\) 为\term[euclid-hanshu]{Euclid 函数}[Euclidean function]。
\end{definition}

本书不额外要求 \(\delta(a)\leq\delta(ab)\)（\(a,b\ne0\)）。有些文献把这个乘法单调性并入定义。两种约定对函数的要求不同，但得到的 Euclid 整环类相同。事实上，给定满足上述除法条件的 \(\delta\)，可对每个 \(a\in R\setminus\{0\}\) 定义
\[
\widetilde\delta(a)\coloneqq\min_{c\in R\setminus\{0\}}\delta(ac).
\]
由于 \(R\) 是整环，取最小值的集合是非空的非负整数集合。若 \(a,b\ne0\)，\(ab\) 的每个非零倍数也是 \(a\) 的非零倍数，故
\(\widetilde\delta(a)\leq\widetilde\delta(ab)\)。又设 \(a\in R\)、\(b\in R\setminus\{0\}\)，取 \(c\ne0\) 使 \(\delta(bc)=\widetilde\delta(b)\)。用 \(\delta\) 作带余除法，得到
\[
a=(bc)q+r=b(cq)+r,
\qquad r=0\ \text{或}\ \delta(r)<\delta(bc).
\]
当 \(r\ne0\) 时，在最小值的定义中取乘子 \(1\)，可知 \(\widetilde\delta(r)\leq\delta(r)<\widetilde\delta(b)\)。因此 \(\widetilde\delta\) 同时满足除法条件与乘法单调性\cite[Theorem 2.1]{ConradRemarksEuclideanDomains}。商和余数通常不唯一，Euclid 函数也未必唯一；对带余除法的存在性证明，还须另行判断是否给出了可执行的求商步骤。

\ProseParagraphBreak
设 \(R\) 为带 Euclid 函数 \(\delta\) 的 Euclid 整环，\(a,b\in R\) 不全为零。从 \(a,b\) 出发，不断以上一次非零余数作除数，直到余数为零；每次所得非零余数的 \(\delta\) 值都严格下降。这个过程称为\term[euclid-suanfa]{Euclid 算法}[Euclidean algorithm]，也称辗转相除法。
\termidx[zhanzhuanxiangchufa]{辗转相除法}

\begin{proposition}[Euclid 算法]\label{b1:prop:euclidean-algorithm}
设 \(R\) 为 Euclid 整环，\(a,b\in R\) 不全为零。依次作带余除法，直到余数为零，则最后一个非零余数 \(d\) 是 \(a,b\) 的最大公因子，并可由回代求得 \(u,v\in R\)，使 \(d=ua+vb\)。
\end{proposition}
\begin{proof}
若 \(b=0\)，取 \(d=a\)、\(u=1\)、\(v=0\)。否则置 \(r_0=a,r_1=b\)，并在 \(r_j\ne0\) 时选取
\[
r_{j-1}=q_jr_j+r_{j+1},
\qquad r_{j+1}=0\ \text{或}\ \delta(r_{j+1})<\delta(r_j).
\]
非负整数 \(\delta(r_1),\delta(r_2),\ldots\) 不可能无限严格下降，所以某一步得到 \(r_{m+1}=0\)、\(r_m\ne0\)。等式表明，一个元素同时整除 \(r_{j-1},r_j\)，当且仅当它同时整除 \(r_j,r_{j+1}\)：两个方向分别使用减法和加法。因此每一对相邻余数的公因子集合相同，最终等于 \(r_m,0\) 的公因子集合，故 \(d=r_m\) 为所求最大公因子。

为明确回代，令 \((u_0,v_0)=(1,0)\)、\((u_1,v_1)=(0,1)\)，并递推
\begin{equation}\label{b1:eq:extended-euclid-recurrence}
u_{j+1}=u_{j-1}-q_ju_j,\qquad
v_{j+1}=v_{j-1}-q_jv_j.
\end{equation}
由 \(r_0=u_0a+v_0b\)、\(r_1=u_1a+v_1b\) 开始归纳，除法等式给出每个 \(r_j=u_ja+v_jb\)。取 \(j=m\) 即得所求关系。
\end{proof}

在求余数的同时按\eqref{b1:eq:extended-euclid-recurrence}更新两组系数，便得到\term[kuozhan-euclid-suanfa]{扩展 Euclid 算法}[extended Euclidean algorithm]。\cref{b1:alg:extended-euclid}同时输出最大公因子与 Bézout 系数。Euclid 整环的定义只保证所需的商和余数存在；要实际运行算法，还须能有效完成环运算和符合下降条件的带余除法。

\begin{algorithm}[htbp]
\caption{扩展 Euclid 算法}
\label{b1:alg:extended-euclid}
\begin{algorithmsteps}{配有 Euclid 函数 \(\delta\) 的 Euclid 整环 \(R\)，及不全为零的 \(a,b\in R\)；其中环运算与符合 \(\delta\) 下降条件的带余除法均可有效计算。}{\(a,b\) 的一个最大公因子 \(d\) 及 \(u,v\in R\)，使 \(d=ua+vb\)。}
\State 置 \((r_0,u_0,v_0)=(a,1,0)\)、\((r_1,u_1,v_1)=(b,0,1)\)，并令 \(j=1\)。
\While{\(r_j\ne0\)}
\State 求 \(q_j,r_{j+1}\)，使 \(r_{j-1}=q_jr_j+r_{j+1}\)，且 \(r_{j+1}=0\) 或 \(\delta(r_{j+1})<\delta(r_j)\)。
\State 置 \(u_{j+1}=u_{j-1}-q_ju_j\)、\(v_{j+1}=v_{j-1}-q_jv_j\)。
\State 令 \(j\leftarrow j+1\)。
\EndWhile
\State \Return \(d=r_{j-1},\ u=u_{j-1},\ v=v_{j-1}\)。
\end{algorithmsteps}
\end{algorithm}

\subsection{整数环及一元多项式环}
\label{b1:subsec:1102:02}

在 \(\mathbb Z\) 中，取 \(\delta(n)=\lvert n\rvert\)。对 \(b\ne0\)，普通整数除法可取 \(0\leq r<\lvert b\rvert\)，故它是 Euclid 函数。以 \(252,198\) 为例，四次除法所取的商依次是 \(q_1=1,q_2=3,q_3=1,q_4=2\)：
\[
252=198+54,\qquad 198=3\cdot54+36,\qquad
54=36+18,\qquad36=2\cdot18.
\]
按\eqref{b1:eq:extended-euclid-recurrence}同步更新系数，得到下表；每行都满足 \(r_j=252u_j+198v_j\)。
\begin{center}
\begin{tabular}{rrrr}
\toprule
\(j\) & \(r_j\) & \(u_j\) & \(v_j\) \\
\midrule
\(0\) & \(252\) & \(1\) & \(0\) \\
\(1\) & \(198\) & \(0\) & \(1\) \\
\(2\) & \(54\) & \(1\) & \(-1\) \\
\(3\) & \(36\) & \(-3\) & \(4\) \\
\(4\) & \(18\) & \(4\) & \(-5\) \\
\(5\) & \(0\) & \(-11\) & \(14\) \\
\bottomrule
\end{tabular}
\end{center}
最后非零余数为 \(18\)；同一结果也可由回代读出：
\begin{equation}\label{b1:eq:integer-bezout-example}
18=54-36=4\cdot54-198=4\cdot252-5\cdot198.
\end{equation}
这里不仅算出最大公因子，也获得了以后解线性同余式所需的整数系数。

\begin{proposition}[多项式带余除法]\label{b1:prop:field-poly-division}
设 \(F\) 为域，\(f,g\in F[x]\)，且 \(g\ne0\)。存在唯一的 \(q,r\in F[x]\)，使
\[
f=gq+r,\qquad r=0\ \text{或}\ \deg r<\deg g.
\]
因此 \(F[x]\) 关于非零多项式上的次数函数为 Euclid 整环。
\end{proposition}
\begin{proof}
若 \(f=0\) 或 \(\deg f<\deg g\)，取 \(q=0,r=f\)。其余情形设 \(\deg f=n\)、\(\deg g=m\leq n\)，首项系数分别为 \(a,b\)。因为 \(F\) 是域且 \(b\ne0\)，可以减去 \(ab^{-1}x^{n-m}g\)，把 \(f\) 的最高次项消去。所得
\[
f_1=f-ab^{-1}x^{n-m}g
\]
为零或次数严格小于 \(n\)。对次数作归纳，写成 \(f_1=gq_1+r\)，其中 \(r\) 满足要求，则
\(q=ab^{-1}x^{n-m}+q_1\) 给出所求分解。次数是非负整数，故消项步骤只能进行有限次。

若又有 \(f=gq'+r'\) 且余数满足相同次数限制，则
\(g(q-q')=r'-r\)。若 \(q-q'\ne0\)，由\cref{b1:prop:polynomial-degree}，非零多项式的次数相加，使左端次数至少为 \(\deg g\)，而右端为零或次数小于 \(\deg g\)，矛盾。因此 \(q=q'\)，再得 \(r=r'\)。
\end{proof}

域假设在首项系数 \(b\) 的求逆处使用。例如在 \(\mathbb Z[x]\) 中，\(x\) 无法除以常数 \(2\) 得到次数小于 \(0\) 的余数：这种余数只能为零，而 \(x\notin2\mathbb Z[x]\)。这不否认某些特殊除数可做除法，例如首一多项式的首项系数就是单位。

\ProseParagraphBreak
在 \(\mathbb Q[x]\) 中取 \(f=x^3-2x+1\)、\(g=x^2-1\)。逐项相减得到
\[
f=xg+(-x+1),\qquad g=(-x-1)(-x+1).
\]
所以首一最大公因子为 \(x-1\)，且 \(x-1=xg-f\)。从最后非零余数 \(-x+1\) 改为首一者，需要同时将 Bézout 系数乘以 \(-1\)。这种归一化不影响生成的主理想。

\subsection{Gauss 整数的例子}
\label{b1:subsec:1102:03}

考虑复数的子环
\[
\mathbb Z[i]=\{\,a+bi\mid a,b\in\mathbb Z\,\},\qquad i^2=-1,
\]
称为\term[gauss-zhengshu]{Gauss 整数环}[ring of Gaussian integers]。它是整环，因为它嵌入 \(\mathbb C\)。定义范数
\[
N(a+bi)\coloneqq a^2+b^2.
\]
由 \(N(z)=z\overline z\) 得 \(N(zw)=N(z)N(w)\)，非零元素的范数为正整数。一个元素是单位当且仅当范数为 \(1\)：必要性由逆元与范数的乘法性得到；充分性由 \(z\overline z=1\) 得到。因此单位恰为 \(\pm1,\pm i\)。
\symidx{\(\mathbb Z[i]\)：Gauss 整数环}

若 \(\alpha\mid\beta\)，则 \(N(\alpha)\mid N(\beta)\)；反向推断一般不成立。即使范数相同，两个元素也未必相伴。例如 \(N(2+i)=N(2-i)=5\)，但
\[
\frac{2-i}{2+i}=\frac{3-4i}{5}\notin\mathbb Z[i],
\]
所以 \(2+i\nmid2-i\)，二者也不相伴。范数为整数素数足以证明一个 Gauss 整数不可约，却不是必要条件。例如 \(3\) 在 \(\mathbb Z[i]\) 中不可约（见\cref{b1:ex:11-gaussian-rational-prime}），但 \(N(3)=9\)。

\begin{proposition}\label{b1:prop:gaussian-euclidean}
设
\[
R=\mathbb Z[i]=\{\,a+bi\mid a,b\in\mathbb Z\,\},\qquad
N(a+bi)=a^2+b^2.
\]
则 Gauss 整数环 \(R\) 关于 \(R\setminus\{0\}\) 上的范数 \(N\) 是 Euclid 整环。
\end{proposition}
\begin{proof}
对 \(\alpha,\beta\in\mathbb Z[i]\)、\(\beta\ne0\)，在复数域中写 \(\alpha/\beta=s+ti\)。分别选最近的整数 \(m,n\)，使
\(\lvert s-m\rvert\leq1/2\)、\(\lvert t-n\rvert\leq1/2\)。置 \(q=m+ni\)、\(r=\alpha-\beta q\)，则
\[
N(r)=N(\beta)\bigl((s-m)^2+(t-n)^2\bigr)
\leq\tfrac12N(\beta)<N(\beta).
\]
因 \(q,r\in\mathbb Z[i]\)，这正是所需带余除法；若 \(r=0\)，也满足定义。
\end{proof}

例如 \(5=(2-i)(2+i)\)，所以 \(2+i\) 整除 \(5\)。又 \(N(2+i)=5\) 为整数素数，任何分解 \(2+i=\alpha\beta\) 都使两个非负整数范数的积为 \(5\)，其中一个必为 \(1\)；故 \(2+i\) 不可约。在 \(\mathbb Z\) 中不可约的 \(5\)，进入 \(\mathbb Z[i]\) 后已经分裂。这提示我们：讨论不可约性之前，工作环必须说清楚。
